\chapter*{前言}
\addcontentsline{toc}{chapter}{前言}

这本书的写作过程比我们预想的要长得多。数字经济这个领域真的太新了，新到每隔几个月，就有一批新的政策、新的案例、新的争论需要补进文稿。改到后来我甚至有点犹豫：是不是应该再等等，等这个学科再"稳定"一点？但我想，考试不会等待任何人。所以它现在和大家见面了，完成先于完美。

数字经济专业硕士从设立到今天不过几年时间，同学们备考时面对的困境是真实且同质的：官方教材和各类报告把概念罗列得很全，但读完之后还是不知道这些概念到底是什么意思、彼此之间是什么关系、考试会怎么考。市面上的资料要么是照抄教材，要么是零散的热点汇编。我们做这本书，就是想填补这个空缺，做一本真正“讲得懂”“看得懂”的数字经济教材。

怎么才能讲得懂看得懂呢？很简单，\textbf{官方教材负责"规定"，这本书负责"解释"。}每一个概念，为什么这么定义、每个限定词意味着什么、边界在哪里，都得讲清楚；每一个能写成模型的结论，都一步步推给你看。没有模型可写的部分，就用案例讲：淘宝怎么用免费策略击溃了易趣，滴滴和快的为什么烧掉那么多钱，P2P 是怎么倒的，数字人民币为什么是现在这个设计。这首先需要我自己对这门科目足够掌握，才能写出更好的教材。

熟悉我们体系的朋友应该知道，我们的课程一直强调两件事：\textbf{第一性原理}和\textbf{概念框架}。数字经济这个学科尤其需要这两样东西。它的表象极其纷繁——平台、数据、算法、数字货币、数字贸易、数据资产入表，看起来互不相干。但拆到第一性，你会发现它们反复在用同一批工具：网络外部性、双边市场、价格歧视、增长核算、货币乘数。这些工具本书都单独成章讲透了，后面每一章都在复用它们。而概念框架的意义在于：数字经济考试里一定会出现你从没见过的名词和案例，这时候框架会帮你判断该调用哪组工具，是竞争分析，还是增长分析，还是制度分析？我们要求你能推理出正确的作答路径，而不是临场编造。这是你战胜竞争对手的基础。

这本书是《一本通关》系列的成员书目，全书以 0258 各校专业课大纲为主线搭建，同时标注了学硕宏微观的关联考点，定位依然是科晟全程班的教材。今年我们做了一个大胆的尝试：让一本教材能够覆盖本科与硕士阶段期末考试、数字经济硕士（专硕、学硕）考研、经济学考研知识补充。因此本书的难度与官方教材相比略有加大，但我们认为这是有必要的。我希望我们的同学们可以不要怕难，认真把这本书学懂，从更深层次地掌握数字经济的基本原理与底层逻辑。

我们承认，书里一定还有不足。数字经济变化太快，教材的修订会是一场没有终点的长跑。如果发现错误或讲不清楚的地方，随时告诉我，下一个版本里它就会被改掉。快速迭代是我们的优势之一。

我们越来越确信一件事：真正好的辅导是把路标擦亮。这本书里的每一个推导、每一个案例，都是我们帮你擦亮的路标。路还是要自己走，但至少不用在黑暗里摸索了。

祝好！

\begin{flushright}
科晟智慧产品研发团队\quad 于北京\\
二〇二六年八月
\end{flushright}
